% ===== §6 The Hermeneutics of Shared Happiness =====
\section{The Hermeneutics of Shared Happiness: Resonance, not Merger}
\label{sec:coexistence}

\noindent
\emph{The place of this chapter in the argument (the seam, marked first, to forestall repetition).} \xref{sub:gen-hermeneutic} established the \emph{single-subject} conferral of meaning (happiness is the meaning I confer on my situation, given in the symbolic); \xref{sub:gen-relational} established the \emph{two-subject} ontological co-existence (relation prior, the individual a knot, the autonomy of the knot preserved). At their intersection lies a question neither has asked: \emph{since happiness is meaning given in the symbolic, a ``shared'' happiness would require some ``shared'' conferral of meaning---so how can two subjects, each cut by its own signifier, share a meaning without one's symbolic order swallowing the other's?} This chapter advances at that intersection. It does \emph{not} re-prove ``the necessity of hermeneutics'' (\xref{sub:gen-hermeneutic} is done); it stands on that result and asks after its two-subject corollary.

\paragraph{The chapter's distinctive thesis (one sentence).} Shared happiness is not two subjects falling into one symbolic order (that is the hermeneutic annexation of one by the other) but the \emph{resonance} of two meaning-worlds that remain translatable to one another without being exhausted by the translation---each conferring meaning in its own signifiers, each open to the other, neither required to surrender its symbolic order. The ``we'' of happiness is the resonance of two meaning-worlds unexhausted by each other, not their \emph{merger}.

\subsection{The question: does shared happiness require shared meaning?}
\label{sub:coex-question}

Relaying \xref{sub:gen-hermeneutic}: happiness is meaning conferred in the symbolic. It follows that two people being happy ``together'' is not a juxtaposition of two private meanings but seems to call for a shared meaning---the happiness of ``watching the rain together'' belongs to the ``we,'' not to two side-by-side ``I''s. But the symbolic order (the order of signifiers, the big Other) is a structured order prior to the individual. Hence the sharpening: is shared meaning a conferring of meaning \emph{within one and the same} symbolic order? \emph{Whose} order? The question is at once one of power and of justice (cf.\ Paper~XVII \citep{paperxvii}).

\subsection{The danger: hermeneutic annexation---the symbolic form of the dissolution of co-existence}
\label{sub:coex-annexation}

The most dangerous ``shared happiness'' is precisely this: one party's symbolic order becomes the shared order, and the other confers meaning \emph{within} it---so that the other's happiness is expressed in \emph{another's signifiers}, is translated, is \emph{aphasic} in its own meaning-world. This is the true form, at the \emph{symbolic-hermeneutic} level, of ``mutual dissolution'': not the phase-capture of the dynamical level (a language lacking the symbolic, treated in the ontological and dynamical registers of \xref{sub:gen-relational} and \xref{sub:gen-twomodes}), but \emph{one meaning-world annexing the other}.

Its concealment (isomorphic with the ``regime bleeding'' of Paper~XVII \citep{paperxvii}): hermeneutic annexation often wears the face of ``how alike we are,'' ``we think the same''---and ``thinking the same'' may be precisely that one party has migrated wholesale into the other's signifiers and lost its own place of conferral. The felt happiness of concord may be the quiet that follows a completed annexation.

The criterion (at the symbolic level, distinct from the criteria of the other registers): in this being-together, does the weaker party \emph{still retain the capacity to confer meaning in its own symbolic order}? Or can its meaning now be expressed only through the stronger party's signifiers? ``Whether subjectivity is preserved'' here takes the concrete form of \emph{whether the place of conferral is preserved}.

\subsection{The positive form: translatable yet not exhausted by translation---the structure of resonance}
\label{sub:coex-resonance}

Healthy co-existence does not abolish the differing symbolic orders; it maintains \emph{translatability} between two meaning-worlds: I can open onto your meaning, can grasp you in your language, \emph{without requiring you to surrender your language}---and conversely. The crucial qualifier, \emph{yet not exhausted by translation}: translation always leaves a remainder (the inter-subjective version of the symbolic--real gap of \xref{sub:gen-drive}---your meaning-world has, for me, a remainder no signifier of mine can exhaust). Healthy co-existence \emph{respects this remainder}, and does not force it flat into ``we are wholly the same.'' Resonance is just this: two meaning-worlds, of different frequency yet answering to one another, rather than the monotone left when one has swallowed the other.

\emph{Connecting to the epistemic hospitality of Paper~XIII \citep{paperxiii}.} Paper~XIII established that ``to present something real in the other's epistemic language'' is a correction of epistemic injustice, a respect rather than a manipulation (provided what is presented is real). Here is its positive employment in the theory of happiness: to open onto the other's symbolic order, to share meaning in a manner the other can grasp, \emph{while retaining one's own place of conferral}---this is neither annexation (the other is not abolished) nor aphasia (one's own is not surrendered). Resonance $=$ a \emph{mutual} epistemic hospitality.

Thus the right reading of ``shared meaning'' is not ``one identical meaning'' but a \emph{continual, remainder-leaving, mutual translation and answering} between two meaning-worlds. The happiness of the ``we'' is generated in this back-and-forth, not in a merger.

\subsection{Taking up the earlier material as subordinate illustration: waiting / shared slowing}
\label{sub:coex-waiting}

The ``waiting / shared slowing'' distinction is here \emph{re-placed} as a concrete illustration of the above symbolic-level thesis (no longer an independent load-bearing point):
\begin{itemize}
  \item pathological traction $=$ unilateral phase-capture (one party's autonomous orientation crushed)---answering, at the dynamical level, to the hermeneutic annexation of \xref{sub:coex-annexation};
  \item healthy traction $=$ bilateral, voluntary mutual yielding, falling into a common rhythm that equals neither party's original pace (``watching the evening together'')---answering, at the dynamical level, to the resonance of \xref{sub:coex-resonance};
  \item the criterion (dynamical level): is the slowing bilateral or unilateral? Is the new rhythm negotiated \emph{between} the two, or captured by the stronger? The difference lies in whether the autonomous term is preserved---the side-image, at the dynamical level, of the symbolic-level criterion (whether the place of conferral is preserved).
\end{itemize}
Connecting to \xref{sub:modes-waiting} (waiting as the proceeding toward a common rhythm) and \xref{sub:neuro-waiting} (the two neural systems of voluntary versus forced waiting): the three layers---symbolic, dynamical, neural---are three layers of one distinction, of which \xref{sub:coex-resonance} is the load-bearing symbolic statement.

\paragraph{An emendation to Paper XVIII (footnote-level, not load-bearing).} On the strength of this chapter, one emendation to Paper~XVIII \citep{paperxvii}: ``shared slowing / voluntary mutual yielding'' is to be marked as a \emph{mechanism of realising} generative balance, not a variant of generative dissolution. What Paper~XVIII refused was \emph{unilateral, condescending} waiting; that refusal is not withdrawn---but this distinction must be added, lest Paper~XVIII read as a prohibition of tenderness. This emendation is a by-product of the chapter, not its thesis; the thesis is the resonance-structure of \xref{sub:coex-resonance}.

\subsection{Close}
\label{sub:coex-close}

Shared happiness is the resonance of two meaning-worlds unexhausted by each other, sustained by a mutual epistemic hospitality, with the mutual preservation of the place of conferral as its criterion of justice. In keeping with the voice-discipline of the paper, this ``resonance'' is likewise not a pinned-down essential form but a saying in motion---what resonance concretely looks like emerges with the two persons' historical and material conditions, and is not foreclosed here.
